\chapter{Örnek Ders: Lineer Cebir}
\section{Genel Yapı ve Kullanım}
Bu bölüm, ders notlarını haftalara ve konulara göre düzenlemek için örnek bir yapı gösterir.

% Set the course for this chapter
\Course{Lineer Cebir (A)}

\section{Hafta 1: Vektörler ve Temel Kavramlar}
\Week{1}

\subsection{Vektör Tanımı}
\Topic{Vektörler}
Bir vektör, büyüklük ve yönü olan bir nesnedir.

\begin{definition}
Bir $n$-boyutlu vektör, $\mathbb{R}^n$ kümesinin bir elemanıdır.
\end{definition}

\begin{theorem}
Her vektörün normu sıfırdan büyüktür veya sıfıra eşittir.
\end{theorem}

\begin{exampleauto}
$\mathbf{v} = (1,2,3)$ bir $\mathbb{R}^3$ vektörüdür. Normu $\|\mathbf{v}\|=\sqrt{1+4+9}=\sqrt{14}$.
\end{exampleauto}

\subsection{Örnek Sorular}
\Topic{Norm ve Uzunluk}
\begin{exampleauto}
Bir vektör $\mathbf{a}=(3,4)$ için normu hesaplayınız.
\end{exampleauto}

\begin{exampleauto}
Bir vektör $\mathbf{b}=(1,-1,2)$ için normu hesaplayınız.
\end{exampleauto}

\section{Hafta 2: Matrisler (Örnek)}
\Week{2}
\Topic{Matris Tanımı}
Matrislerle ilgili kısa bir örnek.

\begin{exampleauto}
$A=\begin{pmatrix}1 & 2\\3 & 4\end{pmatrix}$ matrisinin determinantı $\det(A) = -2$.
\end{exampleauto}

\section{Örnek: Filtrelenmiş Çıktılar}
Aşağıda, tanımlı örnekleri farklı filtrelerle nasıl yazdırabileceğinizi gösteriyorum.

% Print all examples
\subsection{Tüm Örnekler}
\PrintAllExamples

\subsection{Hafta 1'e Ait Örnekler}
\PrintExamplesByWeek{1}

\subsection{"Norm ve Uzunluk" Konusuna Ait Örnekler}
\PrintExamplesByTopic{Norm ve Uzunluk}

